% Punkte und Strecken im Koordinatensystem · Streckenlänge, Mittelpunkt, Teilpunkte · Prüfungs-Fokus Abitur GK – bank/punkte-und-strecken-im-koordinatensystem/e2.jsonl (zusammenbau v0.6, Prüfungs-Fokus)
\einheitenkopf[e2]{Einheit 2 · Streckenlängen, Mittelpunkt und Teilpunkte}
\zweigzeile{Abi GK ×6}
\setcounter{aufgabe}{0}

% Anlauf (ohne Prüfkennung)
\begin{aufgabe}{Streckenlänge, Mittelpunkt, Teilpunkte – Anlauf}
\begin{teile}
\teil Wo muss der Haken an der Stange zwischen $A$ und $B$ hängen, damit er genau in der Mitte ist? Was ist gefragt? \\ \kreuz{eine Länge} \\ \kreuz{ein Punkt} \\ \kreuz{eine Figur}
\teil $A(2 | 1 | 3)$, $B(4 | 4 | 9)$: Länge von $AB$? \leerfeld LE
\teil Ein Zeltmast endet in $S(2 | 2 | 6)$; vier gleich lange Leinen führen zu Heringen am Boden, einer davon steckt in $H(4 | 5 | 0)$ (1 LE = 1 m). Je Leine kommen $40$ cm für die Knoten dazu. Wie viel Leine braucht man insgesamt? \leerfeld[m]
\end{teile}
\end{aufgabe}

\begin{aufgabe}{Streckenlänge, Mittelpunkt, Teilpunkte – Prüfungsaufgaben}
\begin{teile}
\teil Ein Gummiseil ist von $A(3 | 1 | 2)$ nach $B(9 | 9 | 2)$ gespannt (1 LE = 1 m) und dabei um $25\,\%$ länger als ungespannt. Wie lang ist es ungespannt? \hfill (A20) \\ \leerfeld[m]
\teil Ein Seil hängt belastet von $A(1 | 1 | 10)$ bis $B(3 | 7 | 1)$ (1 LE = 1 m) und ist dabei $10\,\%$ länger als unbelastet. Unbelastete Länge? \hfill (A20) \\ \leerfeld[m]
\teil Eine Feder zwischen $A(2 | 3 | 1)$ und $B(3 | 7 | 9)$ (1 LE = 1 cm) ist $50\,\%$ länger als in Ruhe. Länge in Ruhe? \hfill (A20) \\ \leerfeld[cm]
\teil Eine Pyramide ist symmetrisch zur $x_1x_3$-Ebene; $D(1 | 3 | 0)$ ist eine Grundecke, $S(3 | 0 | 6)$ die Spitze. Zeige: $W(2 | 1{,}5 | 3)$ ist der Mittelpunkt der Kante $DS$. Gib den Mittelpunkt $T$ der zu $DS$ symmetrischen Kante an. \hfill (A26) \\ T( \leerfeld | \leerfeld | \leerfeld )
\teil Eine Dachgaube ist symmetrisch zur $x_2x_3$-Ebene; $D(6 | 1 | 0)$ ist eine Ecke, $S(0 | 5 | 8)$ die Firstspitze. Zeige: $W(3 | 3 | 4)$ ist der Mittelpunkt der Kante $DS$. Gib den Mittelpunkt $T$ der zu $DS$ symmetrischen Kante an. \hfill (A26) \\ T( \leerfeld | \leerfeld | \leerfeld )
\teil Ein Kristall ist eine Doppelpyramide, symmetrisch zur $x_1x_2$-Ebene; $D(7 | 4 | 0)$ ist eine Ecke, $S(1 | 2 | 4)$ die obere Spitze. Zeige: $W(4 | 3 | 2)$ ist der Mittelpunkt der Kante $DS$. Gib den Mittelpunkt $T$ der zu $DS$ symmetrischen Kante an. \hfill (A26) \\ T( \leerfeld | \leerfeld | \leerfeld )
\teil Eine Strebe führt von $P(9 | 0 | 3)$ nach $Q(0 | 3 | 0)$; zwei Schellen teilen sie in drei gleich lange Stücke. Koordinaten der Schellen? \hfill (A24)
\teil Ein Kletterseil ist von Haken $A(1 | 4 | 2)$ zu Haken $C(11 | -1 | 7)$ gespannt. Der Karabiner $G$ teilt es so, dass $|AG| : |GC| = 2 : 3$. Koordinaten von $G$? \hfill (A23) \\ G( \leerfeld | \leerfeld | \leerfeld )
\teil Eine Stange reicht von $A(2 | 0 | 8)$ bis $C(6 | 8 | 0)$. Eine Leuchte $G$ sitzt so, dass $|AG| : |GC| = 3 : 1$. Koordinaten von $G$? \hfill (A23) \\ G( \leerfeld | \leerfeld | \leerfeld )
\teil $A(3 | 5 | 2)$ und $B(5 | 6 | 4)$ liegen auf einer Geraden $g$. Gib einen Punkt $D \ne B$ auf $g$ an, der von $A$ so weit entfernt ist wie $B$. \hfill (A20) \\ D( \leerfeld | \leerfeld | \leerfeld )
\teil $A(1 | -2 | 4)$ und $B(4 | 2 | 4)$ liegen auf einer Geraden $g$. Gib einen Punkt $D \ne B$ auf $g$ an, der von $A$ so weit entfernt ist wie $B$. \hfill (A20) \\ D( \leerfeld | \leerfeld | \leerfeld )
\teil $A(0 | 3 | 5)$ und $B(2 | 1 | 6)$ liegen auf einer Geraden $g$. Gib einen Punkt $D \ne B$ auf $g$ an, der von $A$ so weit entfernt ist wie $B$. \hfill (A20) \\ D( \leerfeld | \leerfeld | \leerfeld )
\end{teile}
\end{aufgabe}

\begin{aufgabe}{Streckenlänge, Mittelpunkt, Teilpunkte – Prüfungsaufgaben – weiter}
\begin{teile}
\teil $A(2 | -1 | 2)$; der Punkt $C$ ist festgelegt durch $\overrightarrow{OC} = 3 \cdot \overrightarrow{OA}$. Länge der Strecke $AC$? \hfill (A18) \\ \leerfeld
\teil $A(2 | 3 | 6)$; der Punkt $C$ ist festgelegt durch $\overrightarrow{OC} = - \overrightarrow{OA}$. Länge der Strecke $AC$? \hfill (A18) \\ \leerfeld
\teil $A(4 | 4 | 2)$; der Punkt $C$ ist festgelegt durch $\overrightarrow{OC} = \frac{1}{2} \cdot \overrightarrow{OA}$. Länge der Strecke $AC$? \hfill (A18) \\ \leerfeld
\teil Ein Tunnel verläuft geradlinig von $A(10 | 30 | 3)$ nach $B(40 | 70 | 3)$ (1 LE = 10 m). Ein Zug fährt ab $A$ mit $60$ km/h, gleichzeitig ein zweiter ab $B$ mit $90$ km/h. Welchen Weg im Tunnel legt der Zug ab $B$ zurück, bis sich beide begegnen? \hfill (A19) \\ \leerfeld[m]
\teil Ein Radweg führt geradlinig von $P(2 | 1 | 0)$ nach $Q(8 | 9 | 0)$ (1 LE = 100 m). Anna startet in $P$ mit $15$ km/h, Ben gleichzeitig in $Q$ mit $10$ km/h. Wie weit ist Anna gefahren, wenn sie sich treffen, und wo ist das? \hfill (A19)
\end{teile}
\end{aufgabe}

\medskip
%% Merkkasten Einheit 2: 6 Zeilen, Vorgabe höchstens fünf (3.1) – ungekürzt
\uebersichtskasten{Streckenlänge: die Länge der Strecke AB ist der Betrag des Verbindungsvektors, |AB| = √((b₁ − a₁)² + (b₂ − a₂)² + (b₃ − a₃)²) – erst die Differenzen, dann die Quadrate (nie die Koordinaten selbst quadrieren); das ist der Satz des Pythagoras im Raum. \\ A(1 | 2 | 3), B(5 | 4 | 3): AB = (4 | 2 | 0), |AB| = √(16 + 4 + 0) = √20 ≈ 4,5. \\ Mittelpunkt: M = ½ · (OA + OB) – koordinatenweise der Mittelwert; Nachweis durch Nachrechnen der Formel. \\ M(3 | 3 | 3) ist Mittelpunkt von AB. \\ Teilpunkt: T = A + k · AB liegt für k zwischen null und eins auf der Strecke; Drittelpunkte bei k = 1/3 und k = 2/3. Achtung Teilverhältnis: |AG| : |GB| = 1 : 3 heißt Anteil 1/4, also G = A + 1/4 · AB. \\ Punkt in vorgegebener Höhe: die Strecke als Startpunkt plus k mal Verbindungsvektor ansetzen und k aus der verlangten Koordinate bestimmen.}
